\section{Modifying Measurements to Satisfy Known Conservation Laws}

Let
\[
A=
\begin{bmatrix}
a_1&\cdots&a_m
\end{bmatrix}
\in\mathbb{R}^{n\times m},
\qquad
b=
\begin{bmatrix}
b_1\\ \vdots\\ b_m
\end{bmatrix}
\]
The conservation laws are
\[
A^T x=b
\]
We want the smallest correction \(c\) such that
\[
A^T(y+c)=b
\]
Equivalently,
\[
A^T c=b-A^T y
\]
The minimum-norm solution of this linear equation is
\[
c=A(A^TA)^{-1}(b-A^T y)
\]
The inverse exists because \(a_1,\ldots,a_m\) are independent, so \(A\) has linearly independent columns. Therefore \(A^TA\) is invertible.

The adjusted measurement is
\[
y_{\mathrm{adj}}=y+A(A^TA)^{-1}(b-A^T y)
\]
Now verify the conservation laws:
\[
A^T y_{\mathrm{adj}}
=
A^T y+A^TA(A^TA)^{-1}(b-A^T y)
\]
\[
A^T y_{\mathrm{adj}}
=
A^T y+b-A^T y=b
\]
So the adjusted vector satisfies all known conservation laws exactly.
